\begin{theorem}[The Algorithm 3.1 choice trichotomy]
Let \(\mathcal T\) be a certified Tashkinov series with finite data
\(\mathcal D\), let \(i\) be a coloring index, and suppose
\(\Delta(G)+1\leq k\).  Suppose \(\varphi_i\) is proper, leaves exactly
\(e_0\) uncolored, and \(T_{i+1}\) is closed but not strongly closed
under \(\varphi_i\).  Then exactly the data needed to select the next
Algorithm 3.1 branch are inhabited in one of the following three forms.

First, \(\operatorname{RevisitingExtensionCertificate}(\mathcal D,i)\)
is inhabited.  Or it is not inhabited, and there is an admissible
maximum-defective choice \(Q=(\pi,v,f,u,\delta)\) such that every proper
coloring \(\rho\), leaving exactly \(e_0\) uncolored and stable for
\[
             (T_{i+1},D_i\cup\{\delta\},\pi),
\]
has
\[
 \overline\rho(u)\cap\overline\rho(T_{i+1})=\varnothing.        \tag{SE}
\]
Or no revisiting certificate is inhabited, and there is such a maximum
choice \(Q\), together with a witnessed failure of (SE), a color
\(\gamma\) satisfying
\(\operatorname{ParallelExtensionColor}(\mathcal D,i,\pi,v,\gamma)\),
and a proper coloring \(\sigma\) leaving exactly \(e_0\) uncolored,
stable for
\((T_{i+1},D_i\cup\{\delta\},\pi)\), for which \(v\) is a
\((T_{i+1},\sigma,\{\gamma,\delta\})\)-exit.

Here “witnessed failure” means explicitly that there is a proper
\(\rho\), leaving exactly \(e_0\) uncolored and having the displayed
stability, for which the two displayed missing-color sets are not
disjoint.  The conclusion is this inline disjunction over the existing
RE, maximum-choice, PE-color, stability, and exit predicates.
\end{theorem}
\begin{proof}
Decide first whether
\(\operatorname{RevisitingExtensionCertificate}(\mathcal D,i)\), as
defined by \noderef{RevisitingExtensionCertificate}, is inhabited.  If
it is, return that certificate in the first disjunct.  This gives RE the
priority required by Algorithm 3.1.

Suppose henceforth that no RE certificate is inhabited.  Apply
\noderef{MaximumDefectiveChoiceExists} to the proper coloring
\(\varphi_i\), its unique uncolored edge \(e_0\), and the hypothesis that
\(T_{i+1}\) is closed but not strongly closed.  This gives an admissible
maximum-defective choice
\[
                     Q=(\pi,v,f,u,\delta)
\]
for \((T_{i+1},D_i,\varphi_i)\), including every field of
\noderef{MaximumDefectiveChoice}.

Now decide the exact universal predicate (SE) in the statement.  If it
holds, return \(Q\), the already established nonexistence of RE, and
that universal predicate in the second disjunct.  Notice that the
quantifiers, properness and one-uncolored hypotheses, enlarged stable set
\(D_i\cup\{\delta\}\), and disjointness conclusion are exactly those in
the SE constructor of \noderef{TashkinovSeriesStep}.

It remains to treat failure of (SE).  Negating its nested universal
implications gives a particular partial coloring \(\rho\) such that
\[
 \rho\text{ is proper},\quad
 \rho\text{ leaves exactly }e_0\text{ uncolored},\quad
 \rho\text{ is }(T_{i+1},D_i\cup\{\delta\},\pi)\text{-stable},
                                                                    \tag{1}
\]
and
\[
  \neg\operatorname{Disjoint}
       \bigl(\overline\rho(u),\overline\rho(T_{i+1})\bigr).     \tag{2}
\]
Retain \(\rho\) and (1)--(2), since they constitute the witnessed
negation required in the third disjunct.

By \noderef{ParallelExtensionColorExists}, using
\(\Delta(G)+1\leq k\), choose a color \(\gamma\) satisfying
\noderef{ParallelExtensionColor} for \((\mathcal D,i,\pi,v)\).
In particular,
\[
                         \gamma\in\overline\pi(v).              \tag{3}
\]

We verify the hypotheses of the exit-coloring lemma.  First,
\(T_{i+1}\) is closed under \(\pi\).  Let \(a\) be a boundary edge
whose \(\pi\)-color \(c\) is missing at some \(x\in V(T_{i+1})\).
The vertex clause of \noderef{StableColoring} for the selected coloring
says that \(c\) is also missing at \(x\) under \(\varphi_i\), hence
belongs to \(\overline{\varphi_i}(T_{i+1})\).  A boundary edge is
incident with the tree by \noderef{BoundaryEdges}.  Apply the alternate
stability criterion \noderef{sc1} on its \(\pi(a)\) side; properness and
the unique-uncolored hypotheses hold for both colorings by the maximum
choice.  It follows that
\(\varphi_i(a)=\pi(a)=c\), contradicting the assumed closure of
\(T_{i+1}\) under \(\varphi_i\), namely
\noderef{ColoredVertexSetClosed}.  This proves closure under \(\pi\).

Next, the selected edge \(f\) remains the same exact edge identity.
The fields of \noderef{MaximumDefectiveChoice} say that it is a defective
edge, and \noderef{DefectiveEdge} places it in the exact boundary-color
set; hence \(f\in\partial_G(V(T_{i+1}))\).  They also say
\[
 \operatorname{ends}_G(f)=\{v,u\},\quad
 v\in V(T_{i+1}),\quad u\notin V(T_{i+1}),\quad \pi(f)=\delta.   \tag{4}
\]
Use only symmetry of the unordered pair in (4) to rewrite it as
\(\operatorname{ends}_G(f)=\{u,v\}\), in the outside/inside order
expected below.  No endpoint equality is used to substitute a different
parallel edge.

Finally, (2) is equivalent, for the two finite missing-color sets, to
the existence of a color \(\beta\) in their intersection.  Thus
\[
       \beta\in\overline\rho(u)
       \quad\text{and}\quad
       \beta\in\overline\rho(T_{i+1}).                          \tag{5}
\]
The coloring \(\rho\) and (1), together with (4)--(5), provide the
existential premise of \noderef{LEM_extable}; the maximum choice supplies
properness and unique uncoloredness of \(\pi\), and the preceding
paragraph supplies closure.  Apply that lemma with the color \(\gamma\)
from (3).  It returns a partial coloring \(\sigma\) which is proper,
leaves exactly \(e_0\) uncolored, is
\((T_{i+1},D_i\cup\{\delta\},\pi)\)-stable, and satisfies
\[
                  \operatorname{TreeExit}
                    (T_{i+1},\sigma,\gamma,\delta,v)
\]
in the sense of \noderef{TreeExit}.  Return \(Q\), the retained witness
\(\rho\) of failure of (SE), \(\gamma\) with its prescribed history rule,
and this \(\sigma\) with its exit certificate in the third disjunct.
All three branches preserve the prior decision that RE is unavailable
where required, completing the inline trichotomy.
\end{proof}
